\documentclass{article}
\usepackage{amsmath, amssymb}

\begin{document}

\section*{Pushing--Medium Gravity: Effective Metric from Fermat’s Principle}

This document derives the effective spacetime metric of the Pushing--Medium (PM)
model directly from Fermat’s principle applied to a moving refractive medium.
The derivation uses only the static PM gravitational fields:



\[
\phi(\mathbf r) \equiv \ln n(\mathbf r),
\qquad
\mathbf u(\mathbf r) \equiv \text{gravitational flow field}.
\]



The resulting metric matches the weak-field isotropic form of general
relativity with PPN parameters $\gamma = \beta = 1$.

\section*{1. Fermat’s Principle in a Moving Medium}

For a medium with refractive index $n(\mathbf r)$ and flow velocity
$\mathbf u(\mathbf r)$, the travel time functional for a ray is:



\[
T[\mathbf r(s)]
=
\int
\left[
\frac{n(\mathbf r)}{c}
+
\frac{k_{\mathrm{Fizeau}}}{c^2}\,
\mathbf u(\mathbf r)\cdot\frac{d\mathbf r}{ds}
\right]
ds.
\]



The PM model is empirically consistent with $k_{\mathrm{Fizeau}} = 1$.

The ray path extremizes $T$, yielding the PM ray equations.

\section*{2. Effective Optical Metric}

The optical metric $g^{\mathrm{opt}}_{\mu\nu}$ is defined such that null
geodesics of this metric reproduce the stationary points of $T$. For a moving
medium, the optical metric takes the form:



\[
ds_{\mathrm{opt}}^2
=
-\frac{c^2}{n^2}\,dt^2
+
\left(d\mathbf r - \mathbf u\,dt\right)^2.
\]



Expanding:



\[
ds_{\mathrm{opt}}^2
=
-\frac{c^2}{n^2}\,dt^2
+
d\mathbf r^2
- 2\,\mathbf u\cdot d\mathbf r\,dt
+
|\mathbf u|^2 dt^2.
\]



This metric encodes both the gravitational index field and the flow field.

\section*{3. Weak-Field Expansion}

In the PM gravitational sector:



\[
n = e^\phi \approx 1 + \phi,
\qquad
|\phi| \ll 1,
\qquad
|\mathbf u| \ll c.
\]



Expanding the optical metric to first order:



\[
ds_{\mathrm{opt}}^2
\approx
-(1 - 2\phi)c^2 dt^2
+
d\mathbf r^2
- 2\,\mathbf u\cdot d\mathbf r\,dt.
\]



The $|\mathbf u|^2$ term is second order and is neglected here.

\section*{4. Identification with an Effective Spacetime Metric}

We identify the optical metric with an effective spacetime metric
$g_{\mu\nu}^{\mathrm{eff}}$ whose null geodesics reproduce PM ray trajectories:



\[
ds^2 = g_{\mu\nu}^{\mathrm{eff}} dx^\mu dx^\nu.
\]



Thus:



\[
g_{tt}^{\mathrm{eff}} = -(1 + 2\phi),
\qquad
g_{ti}^{\mathrm{eff}} = -u_i,
\qquad
g_{ij}^{\mathrm{eff}} = (1 - 2\phi)\,\delta_{ij}.
\]



This is the standard weak-field isotropic metric of general relativity with a
gravitomagnetic term $g_{ti}$.

\section*{5. PPN Parameters}

From the effective metric:



\[
g_{tt} = -(1 + 2\phi),
\qquad
g_{rr} = (1 - 2\phi),
\]



the post-Newtonian parameters are:



\[
\gamma = 1,
\qquad
\beta = 1.
\]



These match the general-relativistic values in the weak-field regime.

\section*{6. Interpretation}

\begin{itemize}
    \item The scalar PM field $\phi = \ln n$ plays the role of the Newtonian
          gravitational potential.
    \item The flow field $\mathbf u$ plays the role of a gravitomagnetic vector
          potential.
    \item Fermat’s principle in a moving medium naturally produces the
          weak-field GR metric.
    \item No tensor field is required; the metric emerges from optical
          properties of the PM medium.
\end{itemize}

\section*{7. Summary}

The PM gravitational sector, when interpreted through Fermat’s principle,
produces an effective spacetime metric:



\[
ds^2
=
-(1 + 2\phi)c^2 dt^2
- 2\,\mathbf u\cdot d\mathbf r\,dt
+
(1 - 2\phi)\,d\mathbf r^2,
\]



with PPN parameters $\gamma = \beta = 1$. This provides a direct bridge between
the PM refractive-medium picture and the weak-field limit of general relativity.

\end{document}

